% !TEX root = A0.MiTFG.tex

\chapterbeginx{Conclusiones y líneas futuras}

A partir del conjunto de simulaciones realizadas tanto en el entorno atmosférico como oceánico, pueden extraerse las siguientes conclusiones generales:

\begin{enumerate}
\item \textbf{Predominio de la dispersión cromática.} En todos los escenarios lineales analizados, la dispersión cromática es el mecanismo dominante en el ensanchamiento temporal de los pulsos ópticos. En la atmósfera, su influencia se ve parcialmente compensada por la baja dispersión del aire, mientras que en el medio oceánico, los altos valores de $\beta_2$ característicos del agua de mar producen un ensanchamiento mucho más acusado incluso en trayectos cortos.

\item \textbf{Influencia del \textit{chirp} y del signo del producto $\boldsymbol{\beta_2 C}$.} El signo del \textit{chirp} inicial determina si el pulso se expande o se comprime parcialmente en las primeras etapas de propagación. Cuando $\beta_2 C > 0$, la dispersión y el \textit{chirp} actúan en la misma dirección, amplificando el ensanchamiento. En cambio, si $\beta_2 C < 0$, se produce una compresión inicial del pulso que puede reducir temporalmente la dispersión acumulada. 

\item \textbf{Dependencia con la distancia de propagación.} El ensanchamiento crece de forma no lineal con la distancia, y su aumento se vuelve muy acusado al superar la longitud de dispersión. En el océano, incluso decenas de metros bastan para alcanzar valores de FWHM varias órdenes de magnitud superiores al pulso inicial, confirmando la fuerte dispersión del medio acuático.

\item \textbf{Efecto de la turbulencia.} En la atmósfera, la turbulencia puede adquirir relevancia para valores altos del parámetro estructural ($C_n^2 \ge 10^{-13}\,\mathrm{m^{-2/3}}$), llegando a igualar la contribución de la dispersión en condiciones extremas. En el medio oceánico, sin embargo, los resultados obtenidos con el modelo OTOPS indican que la turbulencia tiene un efecto prácticamente despreciable frente a la dispersión cromática, incluso para gradientes térmicos elevados.

\item \textbf{Efecto de la dispersión de tercer orden ($\boldsymbol{\beta_3}$).} En la atmósfera, el término $\beta_3$ resulta insignificante para pulsos de duración superior a unos pocos cientos de femtosegundos. En el océano, sin embargo, la dispersión de tercer orden introduce ligeras asimetrías en la envolvente temporal y pequeñas oscilaciones en los flancos del pulso, especialmente visibles en simulaciones con pulsos muy cortos o con \textit{chirp} alto. Su contribución sigue siendo secundaria, pero no despreciable.

\item \textbf{Dependencia con la anchura temporal inicial.} Se ha verificado que, a medida que disminuye la FWHM del pulso, su ancho espectral aumenta y la dispersión se intensifica. Pulsos ultracortos $(\text{FWHM} \leq 0.1 \text{ ps})$ presentan factores de ensanchamiento de varios órdenes de magnitud, mientras que pulsos más largos mantienen una propagación mucho más estable. Este comportamiento se ha observado en ambos medios.

\end{enumerate}


El estudio desarrollado en este trabajo se ha centrado en la propagación lineal de pulsos ópticos ultracortos y en la cuantificación del ensanchamiento temporal inducido por la dispersión cromática y la turbulencia, tanto en la atmósfera como en el océano. Sin embargo, los resultados obtenidos abren múltiples líneas de investigación que pueden explorarse en el futuro con el objetivo de profundizar en la compresión del canal óptico y de optimizar las prestaciones de los sistemas de comunicación en medios dispersivos y turbulentos.

En primer lugar, resulta fundamental avanzar hacia el desarrollo de modelos físicos más realistas para el medio oceánico. Actualmente no existe un modelo de turbulencia marina universalmente aceptado; la mayoría de los enfoques adaptan formulaciones atmosféricas derivadas del modelo de Kolmogorov, lo que genera incertidumbres sobre su validez en entornos acuáticos. Por ello, una línea prioritaria consiste en formular modelos espectrales específicos que incorporen la dependencia del índice de refracción con la profundidad, la temperatura y la salinidad, así como la posibilidad de caracterizar canales verticales y \acs{NLOS}~(\acl{NLOS}) de manera más precisa \cite{Zeng}.

En esta misma línea, resulta igualmente necesario avanzar en el desarrollo de estrategias de compensación y adaptación que permitan optimizar el comportamiento del enlace óptico bajo condiciones dinámicas. En este sentido, el uso de transceptores inteligentes capaces de adaptar dinámicamente la potencia, la modulación o la codificación en función de las condiciones del canal se perfila como una de las estrategias más prometedoras para mejorar la fiabilidad y el rendimiento de los sistemas en medios dispersivos y turbulentos \cite{Zeng}. En el ámbito atmosférico, las técnicas de óptica adaptativa constituyen una herramienta clave para mitigar los efectos de la turbulencia y mejorar la calidad del haz en recepción \cite{Trichili}.

Por otro lado, la integración de sistemas híbridos que combinen canales ópticos con radiofrecuencia o acústica constituye una línea de investigación creciente tanto en comunicaciones atmosféricas como submarinas. En el caso de los enlaces atmosféricos, la combinación de sistemas FSO y RF permite aprovechar las ventajas de ambas tecnologías y mejorar la fiabilidad del enlace frente a condiciones meteorológicas adversas, dado que la atenuación afecta de manera diferente a cada canal (la niebla degrada los enlaces FSO, mientras que la lluvia afecta principalmente a los enlaces RF \cite{Nadeem2009}). En este tipo de configuraciones híbridas, la conmutación entre el enlace óptico y el de radio puede realizarse mediante esquemas de conmutación dura o blanda dependiendo de la arquitectura del sistema \cite{Zhang2009}. De forma análoga, en el entorno submarino se han propuesto arquitecturas híbridas óptico-acústicas y óptico-RF capaces de extender el alcance de las comunicaciones mediante la selección adaptativa del canal \cite{Zeng}. Asimismo, el diseño de redes ópticas submarinas basadas en esquemas \acs{MIMO}~(\acl{MIMO}) \cite{Zhang2016,Zhang2015,Akhoundi2016,Jamali2016} o comunicación cooperativa \cite{Ahmad2012} se plantea como una estrategia eficaz para incrementar la cobertura y la capacidad de transmisión, si bien su implementación práctica sigue siendo un desafío pendiente.


Finalmente, en un horizonte más amplio, las comunicaciones FSO se perfilan como una plataforma clave para el desarrollo de enlaces cuánticos de alta seguridad. En los últimos años, la comunicación cuántica por enlaces FSO, y en particular las aplicaciones basadas en QKD, han experimentado un notable crecimiento, ofreciendo un potencial considerable para mejorar los métodos tradicionales de criptografía y avanzar hacia redes ópticas globales más seguras. No obstante, la implementación práctica de sistemas cuánticos por enlaces FSO requiere superar diversos desafíos técnicos, como la sensibilidad de los estados cuánticos a la turbulencia atmosférica.

En conjunto, las líneas de investigación descritas ponen de manifiesto que el estudio de la propagación óptica en medios dispersivos y turbulentos continúa siendo un campo abierto y en constante evolución. La mejora de los modelos físicos del canal, el desarrollo de técnicas de compensación adaptativa y la exploración de arquitecturas híbridas o cuánticas representan pasos decisivos hacia la consolidación de sistemas ópticos más robustos, eficientes y seguros. Avanzar en estas direcciones permitirá no solo optimizar el rendimiento de los enlaces atmosféricos y submarinos actuales, sino también sentar las bases de futuras redes ópticas integradas capaces de operar de forma fiable en condiciones extremas y de soportar aplicaciones de alta velocidad y sensibilidad cuántica.

\chapterend
